\documentclass[a4paper,10pt]{article}
\usepackage[margin=20mm]{geometry}
\usepackage{graphicx,amsmath,booktabs}
\title{EPC2361 Half-Bridge Design Study}
\author{4, 5 and 6 parallel devices per switch position}
\date{\today}
\begin{document}
\maketitle
\section{Design objective and evidence}
Choose the highest feasible switching frequency up to 100 kHz while exceeding
99.5\% three-phase inverter efficiency and keeping the junction below 125 C.
One half bridge contains $2N$ devices. Three identical legs form the 75 V,
18.75 kW inverter; the efficiency loss budget is 94.22 W overall or 31.41 W
per balanced leg. A half bridge's allocated output power is bookkeeping,
not a separately measured DC/DC efficiency.

\textbf{Status: analytical screening, not a qualified hardware design.}
Single-device SPICE characterization and split-gate studies are retained as
separate reports. Their driver and parasitic network do not represent a
4/5/6-device bank. Fixed overlap times and output-capacitance allowance are
still assumptions. Other converter losses are currently zero; PCB, capacitor,
control and driver-quiescent losses must consume the remaining budget.

\section{Fundamental-cycle loss calculation}
For phase voltage modulation index $m$, $V_{phase,pk}=mV_{dc}/2$ and
\[I_{rms}=\frac{P_{3ph}}{3[mV_{dc}/(2\sqrt{2})]\cos\phi},\qquad
 i(\theta)=\sqrt{2}I_{rms}\sin(\theta-\phi).\]
Three sinusoidal voltage references $r_a,r_b,r_c$ use
\[z=-(\max(r_a,r_b,r_c)+\min(r_a,r_b,r_c))/2\]
for centered continuous SVM, or $z=0$ for SPWM. The upper duty is
$d=(1+r_a+z)/2$. Linear limits are $m\le2/\sqrt{3}$ for SVM and $m\le1$
for SPWM. The separate SPWM example uses $m=0.99$, so it produces less
fundamental voltage and needs more current for the same real power.

With $\delta=t_df_s$, conduction per device is integrated as
\[P_{H,cond}=\left\langle(i/N)^2R_{hot}\max(d-\delta,0)\right\rangle,
\quad P_{L,cond}=\left\langle(i/N)^2R_{hot}\max(1-d-\delta,0)\right\rangle.\]
The angle brackets mean a full fundamental-cycle average (7200 samples by
default). This is a carrier-average approximation, without carrier ripple,
deadtime voltage-error compensation or pulse suppression. It is not a
constant-peak-current calculation. At zero deadtime, leg conduction recovers
$I_{rms}^2R_{hot}/N$ exactly within quadrature precision.

Switching pair energy is evaluated at \emph{each} instantaneous $|i|/N$:
\[E_{pair}=\tfrac12 V_{dc}|i|/N\,(t_{on}+t_{off})+(1+k)E_{oss},\quad
 P_{sw,leg}=Nf_s\langle E_{pair}\rangle.\]
One hard pair per leg per carrier period is assumed; current polarity assigns
it to the upper or lower bank. Complementary capacitive heating is allocated
to that same bank as a screening approximation. For nonlinear measured/SPICE
energy curves, average $E(i)$, not $E(\langle i\rangle)$. Do not add Eoss again
if the energy dataset already includes it.

Reverse conduction during deadtime is $2t_df_sV_{rev}|i|/N$ per conducting
device, allocated to the opposite bank by current polarity. GaN has no silicon
body diode or its reverse-recovery charge, but reverse channel conduction
still dissipates energy. Gate supply power is $2NQ_gV_gf_s$ per leg and is
outside transistor cold-plate heat in this model.

\section{Parameters and characterization requirements}
The EPC2361 datasheet gives maximum 25 C $R_{DS(on)}=1$ milliohm,
typical junction-to-top resistance 0.2 K/W and maximum gate charge 34 nC at
the specified 50 V test condition. Hot resistance uses an approximate typical
125 C multiplier of 1.65 on the maximum room-temperature value. Approximate
$E_{oss}=3.2$ microjoule at 75 V and 5 ns on/off overlap are assumptions.
Gate charge at 75 V, dynamic resistance and switching heat need validation.

Characterize individual-device DC resistance versus temperature and current,
then actual parallel-bank double pulses with the common driver impedance,
individual on/off gate resistors, branch inductances and common-source terms.
Cover zero to peak per-device current, voltage, temperature, gate resistance
and layout variations. Record both devices' voltage/gate stress, branch
current sharing and baseline-corrected dissipative switching energy. Check
energy-window and timestep convergence. Reject extrapolation beyond the
validated current range and preserve the dataset's topology/driver metadata.
Existing single-device 20--60 A target sweeps do not cover this entire requirement.

\section{TIM and cold-plate requirements}
\[R_{TIM}=t/(kA)+R_{contact},\qquad
 T_j=T_w+P_{FET,leg}R_{plate-water}+P_{hot}(R_{jc}+R_{TIM}).\]
$P_{hot}$ is the hotter device's \emph{cycle-average} loss times a 1.2 sharing
allowance. The current model estimates mean steady temperature, not
fundamental-cycle or switching-scale temperature ripple. A transient thermal
network and fundamental frequency are needed for a strict instantaneous
125 C limit. Use margin on the mean estimate.

The example 0.5 mm, 17.8 W/(m K), 15 square mm TIM gives 1.873 K/W per
device before contact resistance. Full-package effective contact area is
idealized. A 0.3 K/W plate is shared by \emph{one half bridge}; it is not
per device. For a single common three-phase plate use all three legs' heat.
The electrically source-connected tops require insulating TIM.
\[R_{plate,max}=\frac{125-T_w-P_{hot}(R_{jc}+R_{TIM})}{P_{FET,leg}},\]
\[R_{TIM,max}=\frac{125-T_w-P_{FET,leg}R_{plate}}{P_{hot}}-R_{jc}.\]
These are conditional tradeoffs, not simultaneous independent maxima. Negative
values mean infeasibility. Converted maximum TIM thickness is a mathematical
bound only, not a mechanically suitable pad selection. Spreading, compression,
contact, insulation and coolant warming require separate verification.

\section{Highest-frequency results under current assumptions}
@RESULTS@
\begin{center}\begin{tabular}{llrrr}\toprule
PWM & N & Frequency (kHz) & Efficiency (\%) & Mean $T_j$ (C)\\\midrule
@ROWS@
\bottomrule\end{tabular}\end{center}
Configurations omitted from this table have no passing grid point.
\begin{center}\includegraphics[width=\linewidth]{../../results/design/design.png}\end{center}
\clearpage
\section{Fundamental waveforms and 100 kHz diagnostics}
\includegraphics[width=\linewidth]{../../results/design/fundamental_cycle.png}

The plotted SVM 100 kHz case is a diagnostic even if it fails the pulse screen.
The strict centered-PWM requirement is $d_{min}/f_s>t_d+t_{pulse,min}$.
At SVM $m=1.15$, $d_{min}=(1-\sqrt3 m/2)/2$; with 20 ns deadtime and
10 ns minimum useful pulse, the highest sampled passing frequency is 67 kHz.
This is a controller/timing assumption, not a semiconductor frequency limit.
A lower modulation index, different pulse handling or verified shorter timing
would change it. Changing modulation also changes current at fixed power.

The CSV includes allowed cycle-average switching energy per device:
\[\langle E_{pair}\rangle_{max}=
\frac{P_{budget,leg}-P_{cond}-P_{dead}-P_{gate}-P_{other,leg}}{Nf_s}.\]
This is a useful target for subsequent parallel-bank characterization.

Source inputs and exact results: data/inputs/design.json and
results/design/summary.json. Source device PDF: EPC2361 datasheet revision 2.4,
20 July 2026. Separate reports retain SPICE methods and voltage-stress results.
\end{document}
